\documentclass[11pt]{article}

\usepackage[letterpaper,margin=1in]{geometry}
\usepackage[T1]{fontenc}
\usepackage[utf8]{inputenc}
\usepackage{lmodern}
\usepackage{microtype}
\usepackage{amsmath,amssymb,mathtools,bm}
\usepackage{enumitem}
\usepackage{xcolor}
\usepackage{booktabs}
\usepackage{array}
\usepackage{hyperref}
\usepackage[most]{tcolorbox}
\usepackage{fancyhdr}
\usepackage{titlesec}
\usepackage{parskip}

\definecolor{qftblue}{HTML}{245A8D}
\definecolor{qftgreen}{HTML}{2F6B4F}
\definecolor{qftred}{HTML}{A23B3B}
\definecolor{qftgray}{HTML}{F4F6F8}
\definecolor{qftgold}{HTML}{C08B2C}

\hypersetup{
  colorlinks=true,
  linkcolor=qftblue,
  urlcolor=qftblue,
  citecolor=qftgreen,
  pdftitle={PHYS 662 Lecture 4: Review of Symmetries and Groups},
  pdfauthor={Nima Lashkari}
}

\pagestyle{fancy}
\fancyhf{}
\lhead{PHYS 662: Quantum Field Theory I}
\rhead{Lecture 4}
\cfoot{\thepage}
\setlength{\headheight}{14pt}

\titleformat{\section}{\Large\bfseries\color{qftblue}}{\thesection.}{0.6em}{}
\titleformat{\subsection}{\large\bfseries\color{qftgreen}}{\thesubsection}{0.6em}{}
\titleformat{\subsubsection}{\normalsize\bfseries\color{qftred}}{\thesubsubsection}{0.6em}{}

\setlist[itemize]{leftmargin=2em,itemsep=0.25em,topsep=0.4em}
\setlist[enumerate]{leftmargin=2em,itemsep=0.25em,topsep=0.4em}

\newcommand{\dd}{\mathrm{d}}
\newcommand{\e}{\mathrm{e}}
\newcommand{\ii}{\mathrm{i}}
\newcommand{\cH}{\mathcal{H}}
\newcommand{\cL}{\mathcal{L}}
\newcommand{\cS}{\mathcal{S}}
\newcommand{\R}{\mathbb{R}}
\newcommand{\C}{\mathbb{C}}
\newcommand{\comm}[2]{\left[#1,#2\right]}
\newcommand{\SO}{\mathrm{SO}}
\newcommand{\OO}{\mathrm{O}}
\newcommand{\UU}{\mathrm{U}}
\newcommand{\SU}{\mathrm{SU}}

\newtcolorbox{keyidea}[1][]{
  enhanced,
  breakable,
  colback=qftgray,
  colframe=qftblue,
  boxrule=0.8pt,
  arc=1.5mm,
  title={#1},
  fonttitle=\bfseries,
  coltitle=white
}

\newtcolorbox{remarkbox}[1][]{
  enhanced,
  breakable,
  colback=white,
  colframe=qftgreen,
  boxrule=0.7pt,
  arc=1.5mm,
  title={#1},
  fonttitle=\bfseries,
  coltitle=white
}

\newtcolorbox{warningbox}[1][]{
  enhanced,
  breakable,
  colback=white,
  colframe=qftred,
  boxrule=0.7pt,
  arc=1.5mm,
  title={#1},
  fonttitle=\bfseries,
  coltitle=white
}

\begin{document}
\pagenumbering{gobble}

\begin{titlepage}
\centering
\vspace*{1.2in}
{\scshape\Large PHYS 662 --- Quantum Field Theory I\par}
\vspace{0.35in}
{\Huge\bfseries\color{qftblue} Lecture 4\par}
\vspace{0.15in}
{\LARGE\bfseries Review of Symmetries and Groups\par}
\vspace{0.12in}
{\Large\color{qftgreen} From Phonons to Group Representations\par}
\vspace{0.6in}
{\large Nima Lashkari\par}
{\large Fall 2026\par}
\vfill
\begin{minipage}{0.86\textwidth}
\small
The lecture begins with the harmonic chain and asks what its collective
quanta teach us about particles, interactions, continuum limits, and
effective field theory.  It then develops the language needed to discuss
symmetry precisely: active and passive transformations, scalar, vector, and
tensor fields, groups and Lie groups, and linear representations.  The final
part compares orthogonal, unitary, and indefinite orthogonal groups and
emphasizes that the words ``scalar'' and ``vector'' always refer to a
specified symmetry group.

\medskip
Course page: \url{https://nimalashkari.com/quantum-field-theory-2026-part-i/}
\end{minipage}
\vspace{0.5in}
\end{titlepage}

\pagenumbering{roman}
\tableofcontents
\newpage
\pagenumbering{arabic}

\section*{Roadmap}
\addcontentsline{toc}{section}{Roadmap}

Lecture 3 exhibited a field theory emerging from a microscopic lattice.  Its
long-wavelength phonons obey a continuum wave equation with more symmetry than
the underlying lattice.  That example makes symmetry the natural next
subject: we need precise language for distinguishing an exact microscopic
symmetry, an approximate infrared symmetry, and a change of coordinates.

This lecture develops that language in seven stages:
\begin{enumerate}
  \item review the phonon field, its interactions, and its continuum limit;
  \item interpret the low-energy scalar field as an effective theory with an
  emergent limiting-speed symmetry;
  \item review translations and rotations of Euclidean space;
  \item distinguish active transformations from passive coordinate changes;
  \item define scalar, vector, and tensor fields;
  \item define groups, Lie groups, generators, and representations; and
  \item compare $\OO(N)$, $\SO(N)$, $\UU(N)$, $\SU(N)$, and indefinite
  orthogonal groups.
\end{enumerate}

The central theme is that a transformation of space and a transformation of
field components are related but distinct.  A field carries both an argument,
telling us where it is evaluated, and possibly indices, telling us which
representation it belongs to.  In particular, a field may be a scalar under
spatial rotations while being a vector under an internal symmetry.

\section{Review: from phonons to an effective scalar field}

\subsection{The harmonic chain and its quanta}

Consider $N$ masses $m$ on a periodic one-dimensional lattice of spacing
$a$.  If $q_j$ is the displacement of the $j$th mass from equilibrium and
$p_j=m\dot q_j$, the harmonic Hamiltonian is
\begin{equation}
  H_0=\sum_{j=0}^{N-1}\left[
    \frac{p_j^2}{2m}+\frac K2(q_{j+1}-q_j)^2
  \right],
  \qquad q_{j+N}=q_j.
  \label{eq:harmonicchain}
\end{equation}
The discrete Fourier transform diagonalizes this quadratic many-body problem.
For wave numbers in the first Brillouin zone, the exact dispersion relation is
\begin{equation}
  \omega_k^2=\frac{4K}{m}\sin^2\!\left(\frac{ka}{2}\right).
  \label{eq:latticedispersion}
\end{equation}
At long wavelength, $|k|a\ll1$,
\begin{equation}
  \omega_k=c_s|k|\left[1-\frac{(ka)^2}{24}+\cdots\right],
  \qquad
  c_s=a\sqrt{\frac Km}.
  \label{eq:longdispersion}
\end{equation}

After quantization, each nonzero normal mode is a harmonic oscillator:
\begin{equation}
  \hat H_0
  =E_0+\sum_{k\ne0}\hbar\omega_k\hat a_k^\dagger\hat a_k,
  \qquad
  \comm{\hat a_k}{\hat a_{k'}^\dagger}=\delta_{kk'}.
  \label{eq:phononH}
\end{equation}
The vacuum obeys $\hat a_k\lvert\Omega\rangle=0$, a one-phonon state is
$\hat a_k^\dagger\lvert\Omega\rangle$, and repeated creation operators span
the symmetric Fock space
\begin{equation}
  \mathcal F_s(\cH_1)
  =\bigoplus_{n=0}^{\infty}\operatorname{Sym}(\cH_1^{\otimes n}).
  \label{eq:phononfock}
\end{equation}
The site labels are labels of microscopic normal coordinates.  The phonons are
indistinguishable bosonic quanta of collective modes.

\begin{remarkbox}[A bounded band is not a bounded many-body energy]
Because the Brillouin zone is compact, the energy of a \emph{single} phonon is
bounded:
\begin{equation*}
  0\leq \hbar\omega_k\leq2\hbar\sqrt{K/m}.
\end{equation*}
The total energy is nevertheless unbounded, because arbitrarily many phonons
can occupy the modes.  In the formal continuum theory the momentum range is
extended without bound, so even the one-particle energy $\hbar c_s|k|$ becomes
unbounded.
\end{remarkbox}

\subsection{Interactions and particle-number-changing processes}

The exact solvability of Eq.~\eqref{eq:harmonicchain} follows from its
quadratic form.  A generic bond potential contains anharmonic terms,
\begin{equation}
  V(\Delta q)
  =\frac K2(\Delta q)^2
  +\frac{g_3}{3!}(\Delta q)^3
  +\frac{g_4}{4!}(\Delta q)^4+\cdots,
  \qquad \Delta q=q_{j+1}-q_j.
  \label{eq:anharmonicreview}
\end{equation}
In momentum space a cubic interaction has the schematic structure
\begin{equation}
  H_3\sim
  \sum_{k_1,k_2,k_3}
  F(k_1,k_2,k_3)
  q_{k_1}q_{k_2}q_{k_3}
  \delta^{(G)}_{k_1+k_2+k_3,0},
  \label{eq:cubicmomentum}
\end{equation}
where the lattice form factor $F$ comes from finite differences and the last
factor imposes conservation of crystal momentum modulo a reciprocal lattice
vector.  Expressing each $q_k$ in terms of $a_k$ and $a_k^\dagger$ generates
terms such as
\begin{equation}
  aaa,
  \qquad a^\dagger aa,
  \qquad a^\dagger a^\dagger a,
  \qquad a^\dagger a^\dagger a^\dagger.
  \label{eq:numberchanging}
\end{equation}
They describe decay, merging, and virtual processes.  Thus the separate mode
occupation numbers $N_k=a_k^\dagger a_k$, and their sum
$N=\sum_kN_k$, are conserved by $H_0$ but not by a generic anharmonic
Hamiltonian.

\begin{keyidea}[Particle number is an emergent conservation law]
The microscopic theory conserves energy and crystal momentum, but it need not
conserve phonon number.  Phonons remain useful particles when interactions
produce decay widths much smaller than their energies.  A particle is then a
long-lived spectral feature, not a fundamental constituent of the Hamiltonian.
\end{keyidea}

\subsection{The continuum limit and canonical normalization}

Let $x_j=ja$ and suppose $q_j(t)\simeq q(x_j,t)$ varies slowly on the lattice
scale.  The basic Riemann-sum rules are
\begin{equation}
  \frac{q_{j+1}-q_j}{a}\longrightarrow\partial_xq,
  \qquad
  a\sum_j f(x_j)\longrightarrow\int\dd x\,f(x),
  \qquad
  \sum_j\longrightarrow\frac1a\int\dd x.
  \label{eq:riemannrules}
\end{equation}
Introduce the mass density $\rho=m/a$ and the elastic coefficient $C=Ka$.
Then Eq.~\eqref{eq:harmonicchain} becomes
\begin{equation}
  H_{\mathrm{IR}}
  =\frac12\int\dd x\left[
    \rho(\partial_tq)^2+C(\partial_xq)^2
  \right]+\text{higher-derivative terms}.
  \label{eq:qcontinuumH}
\end{equation}
The canonically normalized field and momentum are
\begin{equation}
  \phi=\sqrt\rho\,q,
  \qquad
  \Pi=\partial_t\phi,
  \qquad
  c_s^2=\frac C\rho=\frac{Ka^2}{m}.
  \label{eq:canonicalphi}
\end{equation}
Therefore
\begin{equation}
  \boxed{
  H_{\mathrm{IR}}
  =\frac12\int\dd x\left[
    \Pi^2+c_s^2(\partial_x\phi)^2
  \right],
  }
  \label{eq:continuumH}
\end{equation}
and the equation of motion is
\begin{equation}
  (\partial_t^2-c_s^2\partial_x^2)\phi(x,t)=0.
  \label{eq:soundwave}
\end{equation}
For a one-dimensional chain, the continuum mode expansion uses $\dd k/(2\pi)$:
\begin{align}
  \hat\phi(x,t)
  &=\int_{|k|<\Lambda}\frac{\dd k}{2\pi}
  \sqrt{\frac{\hbar}{2\omega_k}}
  \left[
    \hat a(k)\e^{\ii(kx-\omega_kt)}
    +\hat a^\dagger(k)\e^{-\ii(kx-\omega_kt)}
  \right],
  \label{eq:phimodeexpansion}
  \\
  \hat\Pi(x,t)
  &=-\ii\int_{|k|<\Lambda}\frac{\dd k}{2\pi}
  \sqrt{\frac{\hbar\omega_k}{2}}
  \left[
    \hat a(k)\e^{\ii(kx-\omega_kt)}
    -\hat a^\dagger(k)\e^{-\ii(kx-\omega_kt)}
  \right].
  \label{eq:Pimodeexpansion}
\end{align}
The three-dimensional scalar-field generalization replaces $\dd k/(2\pi)$
by $\dd^3\bm k/(2\pi)^3$.  This is the formula displayed on the updated
slides.  The one-dimensional chain itself uses $\dd k/(2\pi)$.  If the
microscopic displacement is vector-valued rather than scalar, one must also
include polarization labels.

The scale $\Lambda\sim\pi/a$ is a physical ultraviolet cutoff.  At
$|k|a\ll1$, Eq.~\eqref{eq:soundwave} has a Lorentz symmetry with limiting
speed $c_s$.  The exact sine dispersion in Eq.~\eqref{eq:latticedispersion}
breaks that symmetry at short wavelengths.

\begin{warningbox}[Emergent is not exact]
The solid selects a preferred rest frame, and its microscopic dynamics is not
Lorentz invariant with speed $c_s$.  Lorentz symmetry is an accidental
infrared symmetry of the leading derivative terms.  Higher-derivative
operators remember the lattice and violate it.
\end{warningbox}

\subsection{The scalar field as an effective theory}

The continuum scalar field is an effective description of the long-wavelength
sector.  It is not meant to reproduce the microscopic theory at arbitrary
momentum.  Instead, observables are expanded in powers of $ka$, and the
leading terms are fixed by locality and the infrared symmetries.  The lattice
regulates the ultraviolet by cutting off modes with $|k|\gtrsim\pi/a$.

Phonons are nevertheless physically real: they can be created, absorbed, and
scattered, and they appear as poles or narrow peaks in response functions.  A
finite solid has a discrete normal-mode spectrum.  In a macroscopic sample
the modes become dense, but weak interactions can still leave sharp
quasiparticle peaks.  Thus ``particle'' describes a robust excitation, not
necessarily a microscopic constituent.

This example is also a prototype for particle physics.  A relativistic QFT
may itself be an effective theory below some ultraviolet scale.  Predictivity
does not require knowing every microscopic detail; it requires a controlled
expansion and a clear statement of the regime of validity.

\section{Symmetries of Euclidean space}

Let $X$ denote the space on which a theory is defined and let $\Phi$ denote
all its fields.  A transformation $g$ acts on points and fields.  It is a
symmetry when it preserves the physical structure of interest: for example,
the equations of motion, the action, the Hamiltonian, or the quantum
transition amplitudes.  In a Lagrangian field theory the clean statement is
\begin{equation}
  S[\Phi_g]=S[\Phi].
  \label{eq:actionsymmetry}
\end{equation}
The transformed field $\Phi_g$ must include both the relocation of its
argument and any transformation of its components.

The action of a symmetry on observables can still be nontrivial.  Invariance
does not mean that every state or field configuration is fixed.  It means
that transformed configurations obey the same laws and receive the same
action.

\subsection{The Euclidean metric}

On $\R^d$, the squared length and inner product are
\begin{equation}
  |x|^2=\delta_{ij}x^ix^j,
  \qquad
  x\cdot y=\delta_{ij}x^iy^j.
  \label{eq:euclideanmetric}
\end{equation}
Here and below repeated indices are summed.  The Kronecker delta
$\delta_{ij}$ is the Euclidean metric in Cartesian coordinates.

A linear transformation $x'^i=R^i{}_jx^j$ preserves the metric when
\begin{equation}
  R^T R=I,
  \qquad\text{or equivalently}\qquad
  \delta_{ij}R^i{}_kR^j{}_\ell=\delta_{k\ell}.
  \label{eq:orthogonalcondition}
\end{equation}
The matrices satisfying this condition form the orthogonal group $\OO(d)$.
Taking determinants gives $(\det R)^2=1$, so $\det R=\pm1$.  The subgroup
with $\det R=+1$ is the special orthogonal group $\SO(d)$, the group of proper
rotations.  The other component contains reflections combined with rotations.

In two dimensions,
\begin{equation}
  R(\theta)=
  \begin{pmatrix}
    \cos\theta&-\sin\theta\\
    \sin\theta& \cos\theta
  \end{pmatrix},
  \qquad
  R(\theta_1)R(\theta_2)=R(\theta_1+\theta_2),
  \label{eq:SO2rotation}
\end{equation}
parametrizes $\SO(2)$, not all of $\OO(2)$.  Multiplying such a rotation by a
fixed reflection supplies the second component of $\OO(2)$.

\subsection{Translations and the Euclidean group}

Translations act as $x\mapsto x+a$ and preserve distances between points.
Together, rotations, reflections, and translations act affinely:
\begin{equation}
  x\longmapsto Rx+a,
  \qquad R\in\OO(d),\quad a\in\R^d.
  \label{eq:affinetransformation}
\end{equation}
They form the Euclidean group
\begin{equation}
  E(d)=\R^d\rtimes\OO(d).
  \label{eq:Euclideangroup}
\end{equation}
The semidirect product matters.  If $g_r=(a_r,R_r)$, composition gives
\begin{equation}
  g_2g_1=(a_2+R_2a_1,R_2R_1).
  \label{eq:Euclideangrouplaw}
\end{equation}
Rotations act on translation vectors, so the two subgroups do not simply form
a direct product.

\subsection{Active and passive transformations}

The same geometry can be described in two languages.

\begin{center}
\begin{tabular}{>{\bfseries}p{0.19\textwidth} p{0.32\textwidth} p{0.36\textwidth}}
 & Active transformation & Passive transformation \\
\toprule
Meaning & Move or transform the physical configuration. & Relabel the same
configuration with new coordinates. \\
Coordinates & The coordinate system is fixed. & The coordinate labels change. \\
State & The state or field changes. & The physical state is unchanged. \\
\bottomrule
\end{tabular}
\end{center}

For a translation by $a$ in one dimension, the active transformation of a
scalar profile is
\begin{equation}
  x\longmapsto x+a,
  \qquad
  \phi_a(x)=\phi(x-a).
  \label{eq:activetranslation}
\end{equation}
The inverse appears because the value found at $x$ after the translation came
from $x-a$ before it.  In the passive description, let $x'=x-a$.  Scalarity
means that the numerical value assigned to the same geometric point is
unchanged:
\begin{equation}
  \phi'(x')=\phi(x)=\phi(x'+a).
  \label{eq:passivetranslation}
\end{equation}
The active operation by $a$ and the corresponding passive coordinate change
are inverse descriptions.

The rotation analogue makes the inverse especially visible.  For the active
point transformation $x\mapsto Rx$,
\begin{equation}
  \phi_R(x)=\phi(R^{-1}x),
  \qquad\text{equivalently}\qquad
  \phi_R(Rx)=\phi(x).
  \label{eq:activerotation}
\end{equation}
The inverse passive relabeling is $x'=R^{-1}x$, for which
\begin{equation}
  \phi'(x')=\phi(x)=\phi(Rx').
  \label{eq:passiverotation}
\end{equation}
One may instead choose the passive convention $x'=Rx$; then the correct
relation is $\phi'(x')=\phi(R^{-1}x')$.  Both conventions are valid, but the
definition of $x'$ must be stated before deciding whether $R$ or $R^{-1}$
appears in the argument.

\begin{remarkbox}[Convention used below]
Unless stated otherwise, transformations are written actively.  For a map
$x\mapsto g x$, a scalar transforms as
\begin{equation*}
  \phi_g(x)=\phi(g^{-1}x).
\end{equation*}
This convention prevents the most common sign error in translations and the
most common inverse-matrix error in rotations.
\end{remarkbox}

\subsection{Scalar, vector, and tensor fields}

A scalar field assigns one number to each point.  In the active convention,
the affine point transformation $x\mapsto Rx+a$ gives
\begin{equation}
  \phi_g(x)=\phi\!\left(R^{-1}(x-a)\right).
  \label{eq:activescalar}
\end{equation}

A vector field has both an argument and a vector index.  Consequently, its
active transformation is
\begin{equation}
  V_g^i(x)=R^i{}_j
  V^j\!\left(R^{-1}(x-a)\right).
  \label{eq:vectoractive}
\end{equation}
A rank-three contravariant tensor transforms with one copy of $R$ per index:
\begin{equation}
  T_g^{ijk}(x)
  =R^i{}_{i'}R^j{}_{j'}R^k{}_{k'}
  T^{i'j'k'}\!\left(R^{-1}(x-a)\right).
  \label{eq:tensortransform}
\end{equation}
Covariant indices transform with the inverse matrix.  For an orthogonal
matrix in Cartesian coordinates, $R^{-1}=R^T$, so upper and lower Euclidean
indices are easily identified; this simplification should not obscure the
general rule.

For comparison, the inverse passive coordinates associated with the same
geometric motion are $x'=R^{-1}(x-a)$.  The unchanged scalar and the
components of an unchanged spatial vector then obey
\begin{equation}
  \phi'(x')=\phi(Rx'+a),
  \qquad
  V'^i(x')=(R^{-1})^i{}_jV^j(Rx'+a).
  \label{eq:passivefields}
\end{equation}

\begin{keyidea}[Arguments and components transform separately]
The inverse transformation in the argument tells us where the field value
came from.  The matrix multiplying the components tells us which
representation the field carries.  Scalars have no nontrivial component
matrix; vectors have one; higher tensors have one for each index.
\end{keyidea}

\section{Groups, actions, and representations}

\subsection{The group axioms}

A group $G$ is a set with a product satisfying four axioms:
\begin{enumerate}
  \item \textbf{Closure:} $g_1g_2\in G$ for all $g_1,g_2\in G$.
  \item \textbf{Associativity:}
  $(g_1g_2)g_3=g_1(g_2g_3)$.
  \item \textbf{Identity:} there is an element $e$ such that
  $eg=ge=g$.
  \item \textbf{Inverse:} every $g$ has $g^{-1}$ satisfying
  $g^{-1}g=gg^{-1}=e$.
\end{enumerate}
A group is Abelian if $g_1g_2=g_2g_1$ for every pair and non-Abelian
otherwise.

A group action on a space $X$ is a map $G\times X\to X$, written
$(g,x)\mapsto g\cdot x$, such that
\begin{equation}
  e\cdot x=x,
  \qquad
  g_1\cdot(g_2\cdot x)=(g_1g_2)\cdot x.
  \label{eq:groupaction}
\end{equation}
The group and the space it acts on are conceptually different.  For example,
$\SO(3)$ acts on $\R^3$, but $\SO(3)$ is not $\R^3$.

\subsection{Linear representations}

A linear representation is a homomorphism
\begin{equation}
  \rho:G\longrightarrow GL(V)
  \label{eq:representationmap}
\end{equation}
from group elements to invertible linear maps on a vector space $V$, with
\begin{equation}
  \rho(g_1g_2)=\rho(g_1)\rho(g_2),
  \qquad
  \rho(e)=I,
  \qquad
  \rho(g^{-1})=\rho(g)^{-1}.
  \label{eq:representationproperties}
\end{equation}
Different representations can realize the same abstract group on different
vector spaces.  A representation need not be faithful: distinct group
elements may be represented by the same matrix.

The circle group may be written either as angles modulo $2\pi$ or as phases:
\begin{equation}
  \theta\in\R/(2\pi\mathbb Z),
  \qquad
  z=\e^{\ii\theta}\in\UU(1).
  \label{eq:circleU1}
\end{equation}
The same group has the real two-dimensional representation
$\theta\mapsto R(\theta)$ in Eq.~\eqref{eq:SO2rotation} and complex
one-dimensional representations
\begin{equation}
  \rho_n(\theta)=\e^{\ii n\theta},
  \qquad n\in\mathbb Z.
  \label{eq:U1irreps}
\end{equation}

\subsection{Lie groups and infinitesimal generators}

A Lie group is simultaneously a smooth manifold and a group, with smooth
multiplication and inversion.  A one-parameter subgroup near the identity may
be written
\begin{equation}
  \rho(g(\epsilon))
  =I+\epsilon T+\mathcal O(\epsilon^2)
  =\e^{\epsilon T},
  \label{eq:oneparametergroup}
\end{equation}
where $T$ is an infinitesimal generator.  The generators span the tangent
space to $G$ at the identity.  Their commutators close,
\begin{equation}
  \comm{T_a}{T_b}=f_{ab}{}^cT_c,
  \label{eq:liealgebra}
\end{equation}
and define the Lie algebra.  The constants $f_{ab}{}^c$ are the structure
constants in the chosen basis.

For $\SO(2)$,
\begin{equation}
  J=
  \begin{pmatrix}0&-1\\1&0\end{pmatrix},
  \qquad
  R(\theta)=\e^{\theta J}.
  \label{eq:SO2generator}
\end{equation}
There is only one generator, so the Lie algebra is Abelian.

For $\SO(3)$ choose real antisymmetric generators $(J_i)_{jk}=-\epsilon_{ijk}$.
They obey
\begin{equation}
  \comm{J_i}{J_j}=\epsilon_{ijk}J_k.
  \label{eq:SO3realalgebra}
\end{equation}
With Hermitian quantum-mechanical generators $L_i=-\ii\hbar J_i$, the same
algebra becomes
\begin{equation}
  \comm{L_i}{L_j}=\ii\hbar\epsilon_{ijk}L_k.
  \label{eq:angularmomentumalgebra}
\end{equation}
Rotations in three dimensions are noncommutative.  Euler angles provide a
local parametrization only after an order convention is chosen, for example
\begin{equation}
  R(\alpha,\beta,\gamma)
  =R_z(\alpha)R_y(\beta)R_z(\gamma).
  \label{eq:Eulerangles}
\end{equation}
Because the factors do not commute, changing the order changes the
parametrization and generally changes the resulting rotation.

\begin{warningbox}[A generator convention must be stated]
Real antisymmetric rotation generators have commutators without a factor of
$\ii$.  Hermitian generators used for quantum observables have commutators
with $\ii\hbar$.  Mixing the two conventions produces an apparent sign or
factor-of-$\ii$ contradiction.
\end{warningbox}

\section{Orthogonal and unitary groups}

\subsection{\texorpdfstring{$\OO(N)$}{O(N)} and
\texorpdfstring{$\SO(N)$}{SO(N)}}

The orthogonal group is
\begin{equation}
  \OO(N)=\left\{R\in GL(N,\R)\mid R^TR=I\right\}.
  \label{eq:ONdefinition}
\end{equation}
It preserves the Euclidean norm of a real $N$-component vector:
\begin{equation}
  (R\Phi)^T(R\Phi)=\Phi^TR^TR\Phi=\Phi^T\Phi.
  \label{eq:ONinvariant}
\end{equation}
The determinant separates two components,
\begin{equation}
  \det R=+1\quad\text{or}\quad\det R=-1.
\end{equation}
The identity component is
\begin{equation}
  \SO(N)=\left\{R\in\OO(N)\mid\det R=1\right\}.
  \label{eq:SONdefinition}
\end{equation}
Its generators are real antisymmetric matrices, so
$\dim\SO(N)=N(N-1)/2$.

The same matrices may act as spacetime rotations or as internal
transformations.  If $\Phi^a(x)$ is an $N$-component field and
\begin{equation}
  \Phi'^a(x)=R^a{}_b\Phi^b(x)
  \label{eq:internalON}
\end{equation}
without transforming $x$, then $\OO(N)$ acts on an internal index.  The
quadratic combination $\delta_{ab}\Phi^a\Phi^b$ is invariant.

\subsection{\texorpdfstring{$\UU(N)$}{U(N)} and
\texorpdfstring{$\SU(N)$}{SU(N)}}

For an $N$-component complex vector, the preserved inner product is
$\Phi^\dagger\Phi$.  The unitary group is
\begin{equation}
  \UU(N)=\left\{U\in GL(N,\C)\mid U^\dagger U=I\right\}.
  \label{eq:UNdefinition}
\end{equation}
Indeed,
\begin{equation}
  (U\Phi)^\dagger(U\Phi)
  =\Phi^\dagger U^\dagger U\Phi
  =\Phi^\dagger\Phi.
  \label{eq:UNinvariant}
\end{equation}
Taking determinants gives $|\det U|=1$, so $\det U=\e^{\ii\alpha}$ for some
phase $\alpha$.  The special unitary subgroup is
\begin{equation}
  \SU(N)=\left\{U\in\UU(N)\mid\det U=1\right\}.
  \label{eq:SUNdefinition}
\end{equation}
The real dimensions are $\dim\UU(N)=N^2$ and
$\dim\SU(N)=N^2-1$.

\begin{remarkbox}[Spacetime versus internal symmetry]
An $\OO(d)$ spacetime rotation changes the field argument and rotates any
spatial indices.  An internal $\OO(N)$ or $\UU(N)$ transformation mixes field
components at the same spacetime point.  The same abstract matrix group can
appear in either role; what matters is its action.
\end{remarkbox}

\section{Indefinite orthogonal groups}

The orthogonal-group construction extends from a positive Euclidean metric
to a nondegenerate metric of indefinite signature.  With one positive and
\(d\) negative directions, choose
\begin{equation}
  \eta=\operatorname{diag}(1,-1,\ldots,-1).
  \label{eq:MinkowskiMetric}
\end{equation}
The corresponding indefinite orthogonal group is
\begin{equation}
  \OO(1,d)
  =\left\{
    \Lambda\in GL(d+1,\R)
    \ \middle|\
    \Lambda^T\eta\Lambda=\eta
  \right\}.
  \label{eq:LorentzGroupDefinition}
\end{equation}
It preserves the Minkowski inner product:
\begin{equation}
  (\Lambda v)^T\eta(\Lambda w)
  =v^T\Lambda^T\eta\Lambda w
  =v^T\eta w.
  \label{eq:LorentzInnerProduct}
\end{equation}
In index notation,
\begin{equation}
  \Lambda^\rho{}_\mu\eta_{\rho\sigma}
  \Lambda^\sigma{}_\nu
  =\eta_{\mu\nu}.
  \label{eq:LorentzIndexCondition}
\end{equation}

Some authors write \(\OO(d,1)\) after placing the negative entries first.
The two choices describe isomorphic groups, but the notation and metric
signature should be kept consistent within a calculation.

Taking determinants of Eq.~\eqref{eq:LorentzGroupDefinition} gives
\(\det\Lambda=\pm1\).  The condition \(\det\Lambda=1\) defines
\(\SO(1,d)\), but does not by itself select the component connected to the
identity.  The proper orthochronous group additionally preserves the time
orientation and is denoted
\begin{equation}
  \SO^+(1,d).
  \label{eq:properOrthochronous}
\end{equation}
In spacetime applications this group is called the Lorentz group.  Here it is
used only as an example of the general principle that a matrix group is
defined by the form it preserves; boosts and spacetime kinematics are outside
the scope of this lecture.

\begin{remarkbox}[One pattern behind several classical groups]
Each group is characterized by an invariant form:
\begin{align*}
  \OO(N)&:\quad \Phi^T\Phi,\\
  \UU(N)&:\quad \Phi^\dagger\Phi,\\
  \OO(1,d)&:\quad v^T\eta v.
\end{align*}
The difference between Euclidean rotations, unitary internal transformations,
and indefinite orthogonal transformations is encoded in which bilinear or
Hermitian form is preserved.
\end{remarkbox}

\section{A field can be scalar under one group and vector under another}

Let \(x\in\R^2\) label physical position and let
\begin{equation}
  \Phi(x)=
  \begin{pmatrix}
    \phi_1(x)\\[-2pt]
    \vdots\\[-2pt]
    \phi_N(x)
  \end{pmatrix}
  \label{eq:internalMultiplet}
\end{equation}
be a real \(N\)-component field.  Under an internal orthogonal
transformation,
\begin{equation}
  \Phi'(x)=U\Phi(x),
  \qquad
  U\in\OO(N),
  \label{eq:internalVectorTransformation}
\end{equation}
the point \(x\) is unchanged.  The multiplet is therefore a vector under the
internal group.

Now let \(R\in\SO(2)\) rotate physical space.  If
\begin{equation}
  \Phi_R(x)=\Phi(R^{-1}x),
  \label{eq:spatialScalarMultiplet}
\end{equation}
then each component is a scalar under spatial rotations.  There is no
contradiction: the words \emph{scalar} and \emph{vector} refer to different
group actions.

When both symmetries are present, their actions combine as
\begin{equation}
  \Phi_{(R,U)}(x)=U\Phi(R^{-1}x),
  \qquad
  (R,U)\in
  \SO(2)_{\mathrm{space}}\times\OO(N)_{\mathrm{internal}}.
  \label{eq:combinedSpaceInternalAction}
\end{equation}
The spatial group acts on the argument, while the internal group acts on the
component index.

\begin{keyidea}[Representation labels are relative]
A field is never simply a scalar or a vector without qualification.  It is a
scalar, vector, spinor, or tensor with respect to a specified symmetry group
and a specified representation.  The same field can carry several such
labels at once.
\end{keyidea}

This viewpoint is more precise than imagining an additional physical sphere
attached to each point.  The internal vector space is a space of field
components, not necessarily an extra spatial dimension.  A geometric
realization may sometimes be useful, but it is not part of the definition of
an internal symmetry.

\section{Questions to think about}

\begin{enumerate}
  \item Why is the \(k=0\) lattice mode not an ordinary harmonic oscillator?
  \item Which terms generated by a cubic interaction change phonon number?
  Which quantities remain conserved?
  \item Why must the continuum limit keep quantities such as \(\rho\), \(C\),
  and \(c_s\) fixed?
  \item Which terms in the lattice dispersion break the enhanced continuum
  symmetry?
  \item Why does an inverse transformation appear in the argument of an
  actively transformed scalar field?
  \item Why do the matrices \(R(\theta)\) form \(\SO(2)\), rather than all of
  \(\OO(2)\)?
  \item What is the distinction between a group, a group action, and a linear
  representation?
  \item How do the real antisymmetric generators of \(\SO(3)\) differ from
  Hermitian angular-momentum operators?
  \item How can \(\Phi^a(x)\) be a scalar under spatial rotations and a vector
  under an internal \(\OO(N)\)?
\end{enumerate}

\section*{Summary}
\addcontentsline{toc}{section}{Summary}

The harmonic chain reorganizes microscopic displacement variables into
bosonic phonons.  The harmonic theory conserves free phonon number, while
anharmonic interactions couple momentum modes and permit number-changing
processes.  At wavelengths much longer than the lattice spacing, the chain is
described by a scalar field with approximately linear dispersion and an
enhanced long-wavelength symmetry.  The lattice supplies both the ultraviolet
cutoff and the higher-derivative corrections that break this symmetry.

Euclidean transformations illustrate the structure of symmetry.  Active
transformations move configurations; passive transformations relabel the same
configuration.  A scalar, vector, or tensor field is distinguished by the
representation carried by its components in addition to the transformation
of its argument.

A group specifies composition, a group action specifies how a space is
transformed, and a representation realizes group elements as linear maps.
Lie groups are controlled near the identity by generators and their Lie
algebras.  Orthogonal, unitary, and indefinite orthogonal groups are unified
by the forms they preserve.  Finally, representation labels are relative: a
field can be a scalar under spatial rotations and simultaneously a vector
under an internal symmetry.

\end{document}
